% GENERATED FILE. Edit canonical lessons, then run npm run generate:books.
% canonical-source-hash: fnv1a64:c03fcd912ad61d24
% canonical-lessons: MR-C526-segdi, MR-C526-ota, MR-C526-polpat, MR-C526-malne, MR-C526-kisne

\chapter{Kitchen Tools}
\label{ch:mr-kitchen-tools}
\hlchaptermodality{526}

\begin{chapteropening}
\textbf{By the end of this chapter:} \emph{I can talk about the stove, the counter, kneading and grating.}

Complete the last lesson of chapter 526: I can talk about the stove, the counter, kneading and grating.
\end{chapteropening}

\section[\texorpdfstring{\mr{शेगडी} (śegḍī)}{शेगडी (śegḍī)}]{\mr{शेगडी} (śegḍī) — a stove}

\label{lesson:MR-C526-segdi}



\begin{quote}
Before the new one: say the Marathi for to spill, then the Marathi for to rinse.
\end{quote}

\subsection*{You'll want to know: \mr{शेगडी}}
\textbf{\mr{शेगडी}} — \emph{śegḍī} — \textquotedblleft{}a stove\textquotedblright{}.

\begin{quote}
\mr{आईने} \mr{शेगडीवर} \mr{दूध} \mr{ठेवलं}.
\end{quote}

\begin{quote}
\emph{āīne śegḍīvar dūdh ṭhevlaṁ.}
\end{quote}

\begin{quote}
Mother put the milk on the stove.
\end{quote}

Feminine.

\begin{cousinweb}
\textbf{\mr{शेगडी}} (\emph{śegḍī}) is the word for a cooking stove; its origin is not settled.
\end{cousinweb}

\begin{grammarlens}[title={What you've built}]
One more word for describing what you lived and felt.
\end{grammarlens}

\subsection*{Your turn}
\begin{itemize}
\raggedright
  \item \emph{Say it:} \emph{śegḍī}
  \item \emph{Say it:} \emph{śegḍī}, once more
  \item \emph{Say it:} \emph{āīne śegḍīvar dūdh ṭhevlaṁ.} — Mother put the milk on the stove.
  \item \emph{Recall:} say the Marathi for to spill, then the Marathi for to rinse, then say \emph{śegḍī} again
\end{itemize}

\begin{tcolorbox}[breakable,colback=teal!4,colframe=teal!35!black,title={Before you move on}]
What does \mr{शेगडी} mean? (\textquotedblleft{}A stove\textquotedblright{}.) Say it once more.
\end{tcolorbox}
\section[\texorpdfstring{\mr{ओटा} (oṭā)}{ओटा (oṭā)}]{\mr{ओटा} (oṭā) — a kitchen counter; a raised platform}

\label{lesson:MR-C526-ota}



\begin{quote}
Before the new one: say the Marathi for to rinse, then the Marathi for a stove.
\end{quote}

\subsection*{You'll want to know: \mr{ओटा}}
\textbf{\mr{ओटा}} — \emph{oṭā} — \textquotedblleft{}a kitchen counter; a raised platform\textquotedblright{}.

\begin{quote}
\mr{स्वयंपाकघरातला} \mr{ओटा} \mr{स्वच्छ} \mr{आहे}.
\end{quote}

\begin{quote}
\emph{svayampākgharātlā oṭā svaccha āhe.}
\end{quote}

\begin{quote}
The kitchen counter is clean.
\end{quote}

Masculine.

\begin{cousinweb}
\textbf{\mr{ओटा}} (\emph{oṭā}) is inherited: CDIAL lists Marathi \emph{oṭā}, \textquotedblleft{}a raised mass of earth or bricks\textquotedblright{}, under \emph{*ōṭṭā}, \textquotedblleft{}shelter, screen\textquotedblright{} (CDIAL 2544). The kitchen counter is the same raised platform, moved indoors.
\end{cousinweb}

\begin{grammarlens}[title={What you've built}]
One more word for describing what you lived and felt.
\end{grammarlens}

\subsection*{Your turn}
\begin{itemize}
\raggedright
  \item \emph{Say it:} \emph{oṭā}
  \item \emph{Say it:} \emph{oṭā}, once more
  \item \emph{Say it:} \emph{svayampākgharātlā oṭā svaccha āhe.} — The kitchen counter is clean.
  \item \emph{Recall:} say the Marathi for to rinse, then the Marathi for a stove, then say \emph{oṭā} again
\end{itemize}

\begin{tcolorbox}[breakable,colback=teal!4,colframe=teal!35!black,title={Before you move on}]
What does \mr{ओटा} mean? (\textquotedblleft{}A kitchen counter; a raised platform\textquotedblright{}.) Say it once more.
\end{tcolorbox}
\section[\texorpdfstring{\mr{पोळपाट} (poḷpāṭ)}{पोळपाट (poḷpāṭ)}]{\mr{पोळपाट} (poḷpāṭ) — a rolling board}

\label{lesson:MR-C526-polpat}



\begin{quote}
Before the new one: say the Marathi for a stove, then the Marathi for a kitchen counter.
\end{quote}

\subsection*{You'll want to know: \mr{पोळपाट}}
\textbf{\mr{पोळपाट}} — \emph{poḷpāṭ} — \textquotedblleft{}a rolling board\textquotedblright{}.

\begin{quote}
\mr{आजी} \mr{पोळपाटावर} \mr{पोळ्या} \mr{लाटते}.
\end{quote}

\begin{quote}
\emph{ājī poḷpāṭāvar poḷyā lāṭte.}
\end{quote}

\begin{quote}
Grandmother rolls out chapatis on the rolling board.
\end{quote}

Masculine.

\begin{cousinweb}
\textbf{\mr{पोळपाट}} (\emph{poḷpāṭ}) is a Marathi compound of \mr{पोळी}, \textquotedblleft{}chapati\textquotedblright{}, and \mr{पाट}, \textquotedblleft{}a flat board\textquotedblright{}, inherited from Sanskrit \emph{paṭṭa}, \textquotedblleft{}slab\textquotedblright{} (CDIAL 7699).
\end{cousinweb}

\begin{grammarlens}[title={What you've built}]
One more word for describing what you lived and felt.
\end{grammarlens}

\subsection*{Your turn}
\begin{itemize}
\raggedright
  \item \emph{Say it:} \emph{poḷpāṭ}
  \item \emph{Say it:} \emph{poḷpāṭ}, once more
  \item \emph{Say it:} \emph{ājī poḷpāṭāvar poḷyā lāṭte.} — Grandmother rolls out chapatis on the rolling board.
  \item \emph{Recall:} say the Marathi for a stove, then the Marathi for a kitchen counter, then say \emph{poḷpāṭ} again
\end{itemize}

\begin{tcolorbox}[breakable,colback=teal!4,colframe=teal!35!black,title={Before you move on}]
What does \mr{पोळपाट} mean? (\textquotedblleft{}A rolling board\textquotedblright{}.) Say it once more.
\end{tcolorbox}
\section[\texorpdfstring{\mr{मळणे} (maḷṇe)}{मळणे (maḷṇe)}]{\mr{मळणे} (maḷṇe) — to knead}

\label{lesson:MR-C526-malne}



\begin{quote}
Before the new one: say the Marathi for a kitchen counter, then the Marathi for a rolling board.
\end{quote}

\subsection*{You'll want to know: \mr{मळणे}}
\textbf{\mr{मळणे}} — \emph{maḷṇe} — \textquotedblleft{}to knead\textquotedblright{}.

\begin{quote}
\mr{पोळ्यांसाठी} \mr{आईने} \mr{पीठ} \mr{मळलं}.
\end{quote}

\begin{quote}
\emph{poḷyāṁsāṭhī āīne pīṭh maḷlaṁ.}
\end{quote}

\begin{quote}
Mother kneaded the dough for chapatis.
\end{quote}

Past, with -\mr{ने} (he …): \mr{त्याने} \mr{पीठ} \mr{मळलं}, \emph{tyāne pīṭh maḷlaṁ}.

\begin{cousinweb}
\textbf{\mr{मळणे}} (\emph{maḷṇe}) is inherited: CDIAL lists Marathi \emph{maḷṇẽ}, \textquotedblleft{}to rub\textquotedblright{}, under \emph{*marati}, \textquotedblleft{}crushes, rubs\textquotedblright{} (CDIAL 9870); kneading is rubbing and pressing.
\end{cousinweb}

\begin{grammarlens}[title={What you've built}]
One more word for describing what you lived and felt.
\end{grammarlens}

\subsection*{Your turn}
\begin{itemize}
\raggedright
  \item \emph{Say it:} \emph{maḷṇe}
  \item \emph{Say it:} \emph{maḷṇe}, once more
  \item \emph{Say it:} \emph{poḷyāṁsāṭhī āīne pīṭh maḷlaṁ.} — Mother kneaded the dough for chapatis.
  \item \emph{Recall:} say the Marathi for a kitchen counter, then the Marathi for a rolling board, then say \emph{maḷṇe} again
\end{itemize}

\begin{tcolorbox}[breakable,colback=teal!4,colframe=teal!35!black,title={Before you move on}]
What does \mr{मळणे} mean? (\textquotedblleft{}To knead\textquotedblright{}.) Say it once more.
\end{tcolorbox}
\section[\texorpdfstring{\mr{किसणे} (kisṇe)}{किसणे (kisṇe)}]{\mr{किसणे} (kisṇe) — to grate}

\label{lesson:MR-C526-kisne}



\begin{quote}
Before the new one: say the Marathi for a rolling board, then the Marathi for to knead.
\end{quote}

\subsection*{You'll want to know: \mr{किसणे}}
\textbf{\mr{किसणे}} — \emph{kisṇe} — \textquotedblleft{}to grate\textquotedblright{}.

\begin{quote}
\mr{कोशिंबिरीसाठी} \mr{मीराने} \mr{काकडी} \mr{किसली}.
\end{quote}

\begin{quote}
\emph{kośimbirīsāṭhī mīrāne kākḍī kislī.}
\end{quote}

\begin{quote}
Meera grated a cucumber for the salad.
\end{quote}

Past, with -\mr{ने} (he …): \mr{त्याने} \mr{गाजर} \mr{किसलं}, \emph{tyāne gājar kislaṁ}.

\begin{cousinweb}
\textbf{\mr{किसणे}} (\emph{kisṇe}) is a Marathi verb for grating vegetables; its origin is not settled.
\end{cousinweb}

\begin{grammarlens}[title={What you've built}]
One more word for describing what you lived and felt.
\end{grammarlens}

\subsection*{Your turn}
\begin{itemize}
\raggedright
  \item \emph{Say it:} \emph{kisṇe}
  \item \emph{Say it:} \emph{kisṇe}, once more
  \item \emph{Say it:} \emph{kośimbirīsāṭhī mīrāne kākḍī kislī.} — Meera grated a cucumber for the salad.
  \item \emph{Recall:} say the Marathi for a rolling board, then the Marathi for to knead, then say \emph{kisṇe} again
\end{itemize}

\begin{tcolorbox}[breakable,colback=teal!4,colframe=teal!35!black,title={Before you move on}]
What does \mr{किसणे} mean? (\textquotedblleft{}To grate\textquotedblright{}.) Say it once more.
\end{tcolorbox}
